% chapters/design/architecture.tex
% Sección 4.1: Arquitectura del sistema (~2-3 páginas)
% ============================================================================
%
% GUÍA: Presentar la visión global del sistema. Incluir dos diagramas clave.
%
%   --- DIAGRAMA 1: Arquitectura de componentes ---
%   Mostrar cómo encajan frontend y backend (contexto del TFG coordinado):
%     - Frontend (React — TFG de Alejandro Soler)
%     - Backend API (Django + DRF)
%     - Base de datos (PostgreSQL)
%     - Broker de colas (Redis)
%     - Workers (Celery: OCR, email, tareas periódicas)
%     - Almacenamiento de objetos (MinIO)
%     - Watcher SFTP (monitorización de escáner)
%     - Escáner / impresora (fuente de PDFs)
%
%   Usar una figura:
%   \begin{figure}{fig:architecture}{Diagrama de componentes del \ac{scdee}}
%     \image{}{}{architecture_components}
%   \end{figure}
%
%   TODO: Crear la imagen figures/architecture_components.png (o .pdf)
%         con un diagrama de componentes UML o informal.
%         Herramientas recomendadas: draw.io, Excalidraw, PlantUML.
%
%   --- DIAGRAMA 2: Diagrama de despliegue ---
%   Mostrar los contenedores Docker y sus conexiones de red:
%     - django (API REST, puerto 8000)
%     - postgres (BD, puerto 5432)
%     - redis (broker, puerto 6379)
%     - celery-worker (workers OCR/email)
%     - celery-beat (tareas periódicas)
%     - minio (almacenamiento, puerto 9000)
%     - sftp-watcher (monitorización)
%
%   \begin{figure}{fig:deployment}{Diagrama de despliegue con Docker Compose}
%     \image{}{}{deployment_diagram}
%   \end{figure}
%
%   TODO: Crear la imagen figures/deployment_diagram.png (o .pdf).
%
%   --- CONTENIDO TEXTUAL ---
%   1. Explicar la separación frontend/backend vía API REST.
%   2. Justificar la arquitectura de servicios desacoplados
%      (referencia a RNF de escalabilidad: \ref{rnf:scale:decoupled}).
%   3. Describir la comunicación entre componentes:
%      - HTTP/HTTPS entre frontend y API
%      - TCP entre API y PostgreSQL
%      - Redis como broker para Celery
%      - API S3 entre servicios y MinIO
%      - Filesystem para SFTP watcher
%   4. Mencionar la preparación para despliegue en la nube
%      (\ref{rnf:infra:cloud}).
% ============================================================================

% TODO: Redactar la sección de arquitectura y crear los diagramas.
